\section*{الفصل \ref{ch:b3:lithosphere-thermal-models} --- النماذج الحرارية للغلاف الصخري}

\begin{solution}{exo:b3:lithosphere-thermal-models:1}
$L^2/\kappa = \qty{e12}{s}$، \qty{e14}{s} و\qty{e16}{s}: 0.032، 3.2
و\qty{317}{Myr}. ينسى الكيلومتر درجة حرارته الأصلية خلال عشرات الآلاف من
سنة، أمّا المئة كيلومتر فخلال مئات الملايين.
\end{solution}

\begin{solution}{exo:b3:lithosphere-thermal-models:2}
$u = \qty{e-3}{m}/\qty{3.156e7}{s} = \qty{3.17e-11}{m.s^{-1}}$، إذن
$\mathrm{Pe} = 3.17\times10^{-11}\times\num{3e4}/10^{-6} = 0.95$. الانتقال الأفقي
والتوصيل متساويان في الأهمية: الصخر الصاعد يحمل الحرارة نحو الأعلى
ويزيد انحدار المنحنى الحراري قرب السطح، ولو اعتمدنا على التوصيل وحده
لقلّل من تقدير التدرّج الحراري.
\end{solution}

\begin{solution}{exo:b3:lithosphere-thermal-models:3}
$\sqrt{\kappa t} = \sqrt{\num{0.8e-6}\times5\times\num{3.156e13}} =
\qty{11236}{m}$، $\eta = 10000/(2\times11236) = 0.445$ و$T = 1450\times0.471
= \qty{683}{\celsius}$.
\end{solution}

\begin{solution}{exo:b3:lithosphere-thermal-models:4}
$490/\sqrt1 = 490$، $490/5 = 98$ و$490/10 = \qty{49}{mW.m^{-2}}$. يُرجَّح أنّ المسبار
سيقلّل من تقدير القيمة الأولى: القشرة الفتية متشقّقة ومفتوحة أمام
تداول مياه البحر، التي تحمل جزءًا من الحرارة بعيدًا عبر الينابيع الحارّة.
\end{solution}

\begin{solution}{exo:b3:lithosphere-thermal-models:5}
$z = 2\times1.163\sqrt{\kappa t}$: مع $\sqrt{\kappa t} = 15.89$، 31.78
و\qty{50.25}{km}، $z = 37.0$، 73.9 و\qty{116.9}{km}. تتضاعف الطبقة المبرّدة
من 10 إلى \qty{40}{Myr}؛ وبحلول \qty{100}{Myr} تصبح بالفعل أعمق من
صفيحةٍ سماكتها \qty{95}{km}.
\end{solution}

\begin{solution}{exo:b3:lithosphere-thermal-models:6}
$C$ يتناسب طرديًا مع $\alpha$: $\alpha = \num{3.1e-5}\times350/370.5 =
\qty{2.93e-5}{K^{-1}}$. ويتناسب $C$ أيضًا طرديًا مع $T_m - T_0$ (وشاح عند
$1450\times350/370.5 = \qty{1370}{\celsius}$)، وطرديًا مع $\sqrt\kappa$ (معامل
الانتشار $0.8\times(350/370.5)^2 = \qty{0.71e-6}{m^2.s^{-1}}$): العمق وحده
لا يكفي للفصل بينهما، ولهذا يُستعمل التدفق الحراري أيضًا.
\end{solution}

\begin{solution}{exo:b3:lithosphere-thermal-models:7}
$3.1\times1450/\num{95e3} = \qty{47.3}{mW.m^{-2}}$ و$3.1\times1350/\num{125e3}
= \qty{33.5}{mW.m^{-2}}$. يبلغها \omterm{def:b3:lithosphere-thermal-models:half-space}{نموذج تبريد نصف الفضاء} عند $(490/47.3)^2 =
\qty{107}{Myr}$ و$(490/33.5)^2 = \qty{214}{Myr}$: تحافظ الصفيحة الرقيقة الساخنة على
التدفق الحراري لقاع قديم عند \qty{47}{mW.m^{-2}} ابتداءً من نحو \qty{100}{Myr}
، في حين أن الصفيحة الأسمك والأبرد ستدفعه إلى الانخفاض دون
\qty{40}{mW.m^{-2}}، وهي قيمة نصف الفضاء عند \qty{150}{Myr}.
\end{solution}

\begin{solution}{exo:b3:lithosphere-thermal-models:8}
$A_0h_r = 0.4\times60 = \qty{24}{mW.m^{-2}}$، إذن $A_0 =
\num{24e-3}/\num{e4} = \qty{2.4}{\micro W.m^{-3}}$ و$q_m = 60 - 24 =
\qty{36}{mW.m^{-2}}$. عند \qty{35}{km}، $T = \num{36e-3}\times\num{35e3}/2.5 +
\num{2.4e-6}\times\num{e8}/2.5\times(1 - e^{-3.5}) = 504 + 96\times0.970 =
\qty{597}{\celsius}$.
\end{solution}

\begin{solution}{exo:b3:lithosphere-thermal-models:9}
$h_r = \Delta q_0/\Delta A_0 = \num{18e-3}/\num{1.8e-6} = \qty{10}{km}$،
و$q_m = q_0 - h_rA_0 = 66 - 30 = \qty{36}{mW.m^{-2}}$، لأنّ $h_rA_0 =
\num{e4}\times\qty{3.0e-6}{W.m^{-3}} = \qty{30}{mW.m^{-2}}$.
\end{solution}

\begin{solution}{exo:b3:lithosphere-thermal-models:10}
(أ) $Q(t) = \int_0^tk(T_m - T_0)/\sqrt{\pi\kappa t'}\,\mathrm dt' =
2k(T_m - T_0)\sqrt{t/(\pi\kappa)}$. (ب) الحرارة المفقودة هي المحتوى الحراري
المُزال من العمود، $Q = \rho_mc\int_0^\infty(T_m - T)\,\mathrm dz$،
ويعطي التوازن الأيزوستاتي $(\rho_m - \rho_w)(d - d_r) = \rho_m\alpha\int_0^\infty(T_m -
T)\,\mathrm dz = \alpha Q/c$. (ج) $c = 3.1/(3300\times\num{0.8e-6}) =
\qty{1174}{J.kg^{-1}.K^{-1}}$؛ عند \qty{40}{Myr}، $t = \qty{1.262e15}{s}$
و$Q = 2\times3.1\times1450\times\sqrt{\num{1.262e15}/(\pi\times\num{0.8e-6})} =
\qty{2.015e14}{J.m^{-2}}$؛ $d - d_r = \num{3.1e-5}\times\num{2.015e14}/
(1174\times2270) = \qty{2343}{m}$، وهو يساوي $370.5\sqrt{40}$.
\end{solution}

\begin{solution}{exo:b3:lithosphere-thermal-models:11}
الحدّ $n = 3$ يساوي $3104\times(8/\pi^2)\times\tfrac19e^{-9t/36.2} =
280\,e^{-9t/36.2}$ مترًا: \qty{218}{m} عند \qty{1}{Myr}، \qty{23}{m} عند
\qty{10}{Myr} و\qty{1.9}{m} عند \qty{20}{Myr}. ينخفض إلى \qty{10}{m} عندما
$e^{-9t/36.2} = 10/280$، عند $t = (36.2/9)\ln28 = \qty{13.4}{Myr}$.
\end{solution}

\begin{solution}{exo:b3:lithosphere-thermal-models:12}
(أ) من أجل $z \gg h_r$، $e^{-z/h_r} \to 0$. (ب) $q_m = \qty{24}{mW.m^{-2}}$،
$A_0h_r^2/k = \num{16e-3}\times\num{e4}/2.5 = \qty{64}{K}$، و$z = (1300 -
64)\times2.5/\num{24e-3} = \qty{128.8}{km}$. (ج) $q_m = 2.5\times1236/\num{250e3}
= \qty{12.4}{mW.m^{-2}}$: يجب أن يتقلّص \omterm{def:b3:lithosphere-thermal-models:reduced}{التدفق الحراري المختزَل} إلى نحو النصف، كما هو الحال تحت
الكراتون الذي يسمح جذره السميك والبارد بدخول قليل من الحرارة من الأسفل.
\end{solution}

\begin{solution}{pb:b3:lithosphere-thermal-models:1}
\textbf{1.} $\sqrt{\num{0.8e-6}\times\num{3.156e14}} = \qty{15889}{m}$
و$\sqrt{\num{0.8e-6}\times\num{1.2624e15}} = \qty{31778}{m}$.

\textbf{2.} $\eta = 20000/31778 = 0.6294$ و$20000/63556 = 0.3147$:
$T = 1450\times0.6266 = \qty{909}{\celsius}$ عند \qty{10}{Myr}
و$1450\times0.3437 = \qty{498}{\celsius}$ عند \qty{40}{Myr}.

\textbf{3.} $z = 2\times0.7174\sqrt{\kappa t}$: \qty{22.8}{km}
و\qty{45.6}{km}. يبلغ \qty{60}{km} عندما $\sqrt{\kappa t} = 60/1.4348 =
\qty{41.82}{km}$، عند $t = (\num{41.82e3})^2/\num{0.8e-6} = \qty{2.186e15}{s}$،
\qty{69.3}{Myr}.

\textbf{4.} $q = 3.1\times1450/(\sqrt\pi\times15889) = \qty{159.6}{mW.m^{-2}}$ عند
\qty{10}{Myr} و\qty{79.8}{mW.m^{-2}} عند \qty{40}{Myr}؛ $C_q = 159.6\sqrt{10}
= \qty{505}{mW.m^{-2}.Myr^{1/2}}$.

\textbf{5.} $u = \qty{7.92e-10}{m.s^{-1}}$، $\mathrm{Pe} = 7.92\times10^{-10}
\times\num{95e3}/\num{0.8e-6} = 94$. تنتقل الحرارة أفقيًا مع الصفيحة أسرع
بمائة مرّة من انتشارها أفقيًا، فيبرد كلّ عمود
بالتوصيل الشاقولي وحده.

\textbf{6.} $(\num{95e3})^2/\num{0.8e-6} = \qty{1.13e16}{s}$، \qty{357}{Myr}.
عند \qty{150}{Myr}، $\sqrt{\kappa t} = \qty{61.5}{km}$، والطبقة المبرَّدة، وهي
بضعة أضعاف ذلك، تبلغ عمقًا يقع بعيدًا تحت \qty{95}{km}: إنّ نموذج نصف الفضاء لقاع بهذا القِدَم
يفترض أنّ البرودة قد بلغت أعماقًا يتحرّك فيها الوشاح بالحمل الحراري.

\textbf{7.} $(\rho_m - \rho_w)(d - d_r) = \rho_m\alpha\int(T_m - T)\,\mathrm dz$،
كتلتا عمود الظهرة والعمود القديم متساويتان فوق عمق
التعويض. $T_m - T = T_m\operatorname{erfc}\eta$ يهبط دون
\qty{0.5}{\percent} من $T_m$ عندما $\eta > 2$، فلا يضيف الجزء العميق شيئًا.

\textbf{8.} $C = 2\times3300\times\num{3.1e-5}\times1450/2270\times
\sqrt{\num{0.8e-6}\times\num{3.156e13}/\pi} = 0.13069\times2834.8 =
\qty{370.5}{m.Myr^{-1/2}}$؛ $d = 2500 + 370.5\sqrt{10} = \qty{3672}{m}$
و$2500 + 370.5\sqrt{40} = \qty{4843}{m}$.

\textbf{9.} يعطي النموذج $2500 + 370.5\sqrt{42.5} = \qty{4915}{m}$،
\qty{440}{m} أعمق من الوسيط. يتضمّن الوسيط حمل
الرواسب، والجبال البحرية والهضاب البركانية، وانتفاخات البقع الساخنة،
وخلايا على منحدرات مناطق الكسر والهوامش، وكلّها تجعل القاع
أقلّ عمقًا.

\textbf{10.} $C$ يتناسب طرديًا مع $T_m$: $370.5\times1350/1450 =
\qty{344.9}{m.Myr^{-1/2}}$، قريبة من القيمة المُلائَمة 350.

\textbf{11.} تحمل ماء البحر المتداوِل عبر القشرة الفتيّة المتشقّقة جزءًا
من الحرارة إلى الخارج عبر الينابيع الحارّة بين المجسّات.

\textbf{12.} $\tau = (\num{95e3})^2/(\pi^2\times\num{0.8e-6}) =
\qty{1.143e15}{s} = \qty{36.2}{Myr}$؛ السعة هي $3300\times
\num{3.1e-5}\times1450\times\num{95e3}/4540 = \qty{3104}{m}$.

\textbf{13.} $e^{-40/36.2} = 0.33124$، $e^{-360/36.2}/9 = \num{5.3e-6}$
والحدّ $n = 5$ أصغر أيضًا؛ القوس هو $1 - 0.81057\times0.33125 =
0.73150$، و$d = 2500 + 3104\times0.73150 = \qty{4771}{m}$.

\textbf{14.} بالحدّ $n = 3$، $3104\times0.81057\times\num{5.3e-6} =
\qty{0.013}{m}$: عند \qty{40}{Myr}، $e^{-9t/\tau}$ يساوي $10^{-4}$ من
$e^{-t/\tau}$، ولا يبقى إلّا النمط الأوّل.

\textbf{15.} عند \qty{40}{Myr}، $4843 - 4771 = \qty{73}{m}$؛ عند
\qty{50}{Myr}، $2500 + 370.5\sqrt{50} = 5120$ مقابل $2500 + 3104\times(1 -
0.81057\times0.25127) = 4972$: \qty{148}{m}. بالاستيفاء الخطّي، $40 + 10\times
(100 - 73)/(148 - 73) = \qty{43.6}{Myr}$، أي نحو \qty{44}{Myr} (المتسلسلة الكاملة
تعطي 44.0).

\textbf{16.} $q = (3.1\times1450/\num{95e3})[1 + 2(0.33124 + 0.01204 +
0.00005)] = 47.32\times1.68666 = \qty{79.8}{mW.m^{-2}}$، وهي نفسها
لنموذج نصف الفضاء: عند \qty{40}{Myr} لم تظهر القاعدة أثرها بعد
عند السطح.

\textbf{17.} $2500 + 3104 = \qty{5604}{m}$ و$kT_m/L = \qty{47.3}{mW.m^{-2}}$.

\textbf{18.} $d(120) = 2500 + aL\bigl[1 - (8/\pi^2)e^{-120/\tau}\bigr]$، حيث
$a = \rho_m\alpha T_m/\bigl(2(\rho_m - \rho_w)\bigr) = \qty{32.673}{m/km}$
و$\tau = L^2/(\pi^2\kappa)$.

\textbf{19.} $L = 3200/\bigl(32.673\,[1 - 0.81057e^{-120/\tau}]\bigr)$.
من $L = \qty{100}{km}$: $\tau = \qty{40.13}{Myr}$، القوس 0.95925، $L =
\qty{102.10}{km}$؛ ثمّ $\tau = 41.83$، $L = 102.67$؛ $\tau = 42.30$، $L =
102.82$؛ $\tau = 42.43$، $L = 102.87$؛ ويستقرّ عند $L = \qty{102.9}{km}$.

\textbf{20.} $\tau = \qty{42.48}{Myr}$، $aL = \qty{3362}{m}$: $d(40) = 2500 +
3362\times(1 - 0.81057e^{-40/42.48}) = \qty{4799}{m}$، مقابل \qty{4843}{m}
لنموذج نصف الفضاء و\qty{4800}{m} المقاس: تتوافق الصفيحة مع العمرين معًا.

\textbf{21.} $\tau = \qty{42.5}{Myr}$؛ العمق المقارب $2500 + 3362 =
\qty{5862}{m}$.

\textbf{22.} $3362\times0.81057\,e^{-t/42.48} = 100$ يعطي $t = 42.48\ln27.25 =
\qty{140}{Myr}$.

\textbf{23.} الصفيحة التي تلائم قاع المحيط القديم المتخيَّل
يبلغ سمكها \textbf{\qty[text-series-to-math=true]{102.9}{km}}.
\end{solution}
